\section{Algorithm 1 in the present example}
\label{sec:algorithm}

For a homogeneous polynomial $q$ of degree $k$,
\[
\deg(\LieL q)=k,
\qquad
\deg(\LieQ q)=k+1.
\]
Accordingly, each homogeneous dictionary is constructed in two stages:

\[
\boxed{
\text{quadratic generation of degree }k+1
\quad\longrightarrow\quad
\text{linear closure within degree }k+1.
}
\]

For a set $\mathcal Q$ of degree-$k$ commutative monomials, define its linear
closure $\ClL(\mathcal Q)$ as the smallest degree-$k$ set containing
$\mathcal Q$ and all monomials generated by repeated application of
$\LieL$. In this SEIR model, linear closure replaces infectious factors by
exposed factors with the same city index, while preserving total degree and all
node indices.

\begin{flushleft}
\refstepcounter{algorithm}\label{alg:three_city_algorithm}%
\begin{tabularx}{\textwidth}{@{}>{\raggedleft\arraybackslash}p{0.06\textwidth}X@{}}
\toprule
\multicolumn{2}{@{}l@{}}{\textbf{Algorithm \thealgorithm} (Linearly closed graph-induced lifting for the three-city model)}\\
\midrule
1 & Decompose the vector field into $f=f_{\mathrm L}+f_{\mathrm Q}$.\\
2 & Create one quadratic transmission seed $S_iI_j$ for every channel $(i,j)\in\mathcalE$.\\
3 & Apply linear closure to the quadratic seeds. This adds $S_iE_j$ for every retained $S_iI_j$.\\
4 & Differentiate every degree-$k$ coordinate with $f_{\mathrm Q}$ and collect all distinct degree-$(k+1)$ monomials.\\
5 & Canonicalize products as commutative monomials, retain repeated factors, and merge duplicates. Coefficients are combined only within the same scalar derivative equation.\\
6 & Apply iterative linear closure to the new seed set until no same-degree monomial is missing.\\
7 & Choose a deterministic ordering, build the observable vector, and assemble the retained block operator.\\
8 & Verify linear closure. At truncation order $d$, place only the degree-$(d+1)$ terms generated by $f_{\mathrm Q}$ from the top retained block in the residual.\\
\bottomrule
\end{tabularx}
\end{flushleft}

The remainder of the document applies these steps explicitly up to degree
three and then performs a forward check at degree four.

\stepbox{Reading guide}{
The symbols with a hat, such as $\seedM_3^{\mathcalE}$, denote the seed set
created by the quadratic field. The unhatted set $\mathcal M_3^{\mathcalE}$
denotes the final linearly closed dictionary. This distinction is essential:
the seed set records where the nonlinear transmission mechanism first creates
new monomials, while the closure set records every same-degree coordinate
needed for exact propagation under the linear SEIR transitions.}
